\documentclass[11pt,a4paper]{article}
\usepackage[utf8]{inputenc}
\usepackage[T5]{fontenc}
\usepackage[vietnamese]{babel}
\usepackage[a4paper,margin=1.35cm,headsep=4mm,footskip=7mm]{geometry}
\usepackage{tabularx,array,enumitem,fancyhdr,lastpage,pifont,amsmath,amssymb,multicol,longtable}
\setlength{\parindent}{0pt}
\setlength{\parskip}{1.5pt}
\setlength{\headheight}{15pt}
\setlength{\columnsep}{6mm}
\raggedbottom
\pagestyle{fancy}
\fancyhf{}
\lhead{\footnotesize TRUNG TÂM TOÁN TITAN THỦ ĐỨC}
\rhead{\footnotesize PHIẾU BÀI TẬP - LỚP 2 - BUỔI 1}
\cfoot{\footnotesize Trang \thepage/\pageref{LastPage}}
\setlist[itemize]{leftmargin=6mm,itemsep=0.5mm,topsep=0.5mm,parsep=0pt}
\setlist[enumerate]{leftmargin=7mm,itemsep=0.5mm,topsep=0.5mm,parsep=0pt}
\newcommand{\lessonpart}[1]{\vspace{1.5mm}\noindent\textbf{#1}\par\vspace{0.5mm}}
\newcommand{\answerline}{\noindent\dotfill\par\vspace{0.8mm}}
\newcommand{\shortblank}[1]{\makebox[#1]{\dotfill}}
\renewcommand{\arraystretch}{1.25}
\begin{document}
\begin{center}
{\Large\textbf{PHIẾU BÀI TẬP - CHƯƠNG TRÌNH ƯƠM MẦM VÀ PHÁT TRIỂN TÀI NĂNG TOÁN}}\\[1mm]
{\large\textbf{LỚP 2 - BUỔI 1: KHỞI ĐỘNG TƯ DUY TITAN}}\\[-0.5mm]
{\normalsize\textbf{THINK LIKE A MATHEMATICIAN - SUY NGHĨ NHƯ MỘT NHÀ TOÁN HỌC}}
\end{center}
\noindent\textbf{Họ và tên:} \shortblank{70mm} \hfill \textbf{Nhóm:} \shortblank{25mm} \hfill \textbf{Ngày:} \shortblank{30mm}

\textit{Nhắc em: Đừng dừng lại khi đã có một đáp án. Hãy hỏi: Có cách khác không? Làm sao em biết?}

\section*{PHẦN 1. QUAN SÁT VÀ TÌM QUY LUẬT (Notice \& Pattern)}
\textbf{Bài 1.} Điền số còn thiếu: \quad 2, 4, 6, $\ldots$, 10 \qquad và \qquad 3, 6, 9, 12, $\ldots$\\
Em nhận thấy (\textit{I notice}): \answerline
\textbf{Bài 2.} Dãy hình: $\triangle,\ \circ,\ \square,\ \triangle,\ \circ,\ \square,\ldots$ Hình thứ 11 là hình gì? Em có cách tìm mà không cần vẽ đủ 11 hình không?\\
\answerline
\answerline
\textbf{Bài 3.} Có bốn số 6, 8, 10, 13. Chọn một số ``khác nhóm''. Tìm \textbf{ít nhất hai câu trả lời khác nhau} và nêu lí do.\\
\textit{Em nghĩ số ... khác nhóm vì ...}\\
\answerline
\answerline
\textbf{Bài 4.} \textbf{SOLVER $\rightarrow$ CREATOR:} Em tự tạo một quy luật gồm ít nhất 5 số hoặc 5 hình để bạn đoán phần tiếp theo.\\
\answerline

\section*{PHẦN 2. TÌM NHIỀU CÁCH (Number Maker)}
\textbf{Bài 5.} Cho các thẻ số 2, 5, 8 (mỗi thẻ có thể dùng nhiều lần, chỉ dùng phép cộng, trừ). Viết các phép tính có kết quả bằng \textbf{10}. Tìm càng nhiều cách càng tốt.\\
\answerline
\answerline
\textbf{Bài 6.} Điền ba số \textbf{khác nhau} (từ 1 đến 9) vào $\square+\square+\square=12$. Tìm \textbf{ít nhất 4 cách}. Hai cách chỉ đổi vị trí được tính là một cách.\\
\answerline
\answerline
\textbf{Bài 7. (Open Problem)} Lan có một số viên bi. Minh cho Lan thêm một số viên. Bây giờ Lan có 12 viên. Ban đầu Lan có thể có bao nhiêu viên? Viết các khả năng em tìm được:\\
\answerline
Làm sao em biết đã tìm đủ? \answerline

\section*{PHẦN 3. HÌNH VÀ SUY LUẬN (Shape \& Logic)}
\textbf{Bài 8.} Dùng đúng 6 que tính, xếp và vẽ lại các hình \textbf{khác nhau} mà em tạo được. Nếu chỉ xoay hình thì có tính là hình mới không? Vì sao?
\vspace{2.2cm}

\textbf{Bài 9.} Mai, Nam, An đứng thành một hàng. Mai không đứng đầu. Nam không đứng cuối. Liệt kê \textbf{tất cả} cách xếp có thể (viết theo thứ tự từ đầu hàng đến cuối hàng).\\
\answerline
\textbf{Bài 10*. (Strategy Challenge)} Có 10 que tính. Hai người lần lượt lấy 1 hoặc 2 que. Người lấy que cuối cùng thắng. Nếu em đi trước, em lấy mấy que ở lượt đầu để có chiến lược thắng? Ghi lại ít nhất 2 ván chơi và nêu quy luật em tìm được.\\
\answerline
\answerline
\answerline

\section*{TỔNG KẾT CUỐI BUỔI}
\textbf{Khi gặp bài khó, bước tiếp theo của em là gì?} \quad $\square$ Bỏ qua \quad $\square$ Đoán \quad $\square$ Thử một chiến lược khác và kiểm tra.\\
\textbf{Tự đánh giá:} $\square$ Em làm được và giải thích được. \quad $\square$ Em cần một chút hỗ trợ. \quad $\square$ Em cần luyện thêm.

\begin{center}
\fbox{\parbox{0.91\linewidth}{\centering\textbf{5 BƯỚC TITAN:} OBSERVE $\rightarrow$ THINK $\rightarrow$ TRY $\rightarrow$ CHECK $\rightarrow$ EXPLAIN}}
\end{center}
\end{document}
